% =============================================================
%  08_stringhe_io.tex
%  Capitolo 8 -- Stringhe, I/O e parsing
% =============================================================

\chapter{Stringhe, I/O e parsing}

\begin{intuizione}
Tutto quello che hai scritto finora lavora in memoria,
su dati costruiti direttamente nel codice.
Prima o poi i dati vengono da fuori: un file, una riga
digitata al terminale, una stringa da interpretare.

\medskip
Questo capitolo costruisce la catena completa:
dalla sequenza di caratteri grezzi fino a un mini-interprete
che valuta espressioni Lisp.
Ogni sezione aggiunge un livello di astrazione al precedente.
\end{intuizione}

\begin{tcolorbox}[
  enhanced,
  colback=black!2, colframe=black!20,
  arc=3pt, boxrule=0.4pt,
  left=12pt, right=12pt, top=9pt, bottom=9pt,
  before skip=8pt, after skip=10pt
]
\small\color{black!65}

\textbf{La catena di astrazione di questo capitolo.}

\medskip
\begin{enumerate}
  \item \textbf{Sequence.} Il tipo astratto comune a liste, vettori e stringhe.
  \item \textbf{Caratteri e stringhe.} Le operazioni fondamentali.
  \item \textbf{Stream e file.} Leggere e scrivere dall'esterno.
  \item \textbf{\fn{read} e \fn{write}.} Il confine tra codice e dati in Lisp.
  \item \textbf{Tokenizzazione.} Spezzare una stringa in parti significative.
  \item \textbf{Parser RPN.} Prima applicazione: valutare espressioni postfisse.
  \item \textbf{Mini-evaluator.} Valutare un sottoinsieme di S-expression.
\end{enumerate}
\end{tcolorbox}

\vspace{4pt}

% ================================================================
\section{Sequence: il tipo astratto comune}
% ================================================================

In Common Lisp le liste, i vettori e le stringhe sono tutti
sottotipi del tipo astratto \texttt{sequence}.
Questo significa che molte funzioni operano indifferentemente
su tutti e tre.

\begin{concetto}{Sequence}
\keyline{\fn{length}}{numero di elementi di una sequenza qualsiasi.}
\keyline{\fn{elt}}{accesso per indice: \fn{(elt seq i)}.}
\keyline{\fn{subseq}}{sottosequenza: \fn{(subseq seq start end)}.}
\keyline{\fn{map}}{trasforma: \fn{(map type fn seq...)}. Nota: diversa da \fn{mapcar}.}
\keyline{\fn{remove}}{filtra per valore o predicato.}
\keyline{\fn{concatenate}}{concatena: \fn{(concatenate type seq1 seq2 ...)}.}

\boxsep
\detail{\fn{map} e \fn{concatenate} richiedono il tipo del risultato
come primo argomento: \fn{'list}, \fn{'vector} o \fn{'string}.
Questo è diverso da \fn{mapcar}, che restituisce sempre una lista.}
\end{concetto}

\begin{esempio}[\fn{map} e \fn{concatenate} su tipi diversi]
\begin{lstlisting}
;; map su lista -> lista
(map 'list #'1+ '(1 2 3))        ; => (2 3 4)

;; map su stringa -> lista di caratteri trasformati
(map 'list #'char-upcase "ciao") ; => (#\C #\I #\A #\O)

;; map su stringa -> stringa
(map 'string #'char-upcase "ciao") ; => "CIAO"

;; concatenate costruisce il tipo indicato
(concatenate 'string "foo" "bar")  ; => "foobar"
(concatenate 'list '(1 2) '(3 4))  ; => (1 2 3 4)

;; subseq funziona su qualsiasi sequenza
(subseq "abcdef" 2 5)   ; => "cde"
(subseq '(a b c d) 1 3) ; => (B C)
\end{lstlisting}
\end{esempio}

\begin{attenzione}
\fn{map} (su sequenze) e \fn{mapcar} (su liste) non sono la stessa cosa.
\fn{map 'list} è equivalente a \fn{mapcar} solo quando l'input
è una lista. Su stringhe e vettori usare \fn{map}.
Usare \fn{mapcar} su una stringa produce un errore di tipo.
\end{attenzione}

\begin{workbook}
\textbf{Esercizio 136}\\[4pt]
Prende confidenza con le funzioni di sequenza su stringhe
e vettori prima di entrare nelle sezioni più specifiche.
\end{workbook}

% ================================================================
\section{Caratteri e stringhe}
% ================================================================

In Common Lisp i caratteri sono un tipo distinto dalle stringhe.
Una stringa è un vettore di caratteri.

\argomento{Caratteri}

\begin{concetto}{Il tipo \fnt{character}}
\keyline{\fn{\#\textbackslash{}a}}{sintassi letterale per il carattere \texttt{a}.}
\keyline{\fn{char-code}}{codice Unicode del carattere.}
\keyline{\fn{code-char}}{carattere dal codice numerico.}
\keyline{\fn{char-upcase}/\fn{char-downcase}}{conversione maiuscolo/minuscolo.}
\keyline{\fn{char=}, \fn{char<}}{confronto (case-sensitive).}
\keyline{\fn{char-equal}}{confronto case-insensitive.}

\boxsep
\detail{Qualche letterale utile:
\fn{\#\textbackslash{}Space} (spazio), \fn{\#\textbackslash{}Newline} (a-capo),
\fn{\#\textbackslash{}Tab} (tabulazione).}
\end{concetto}

\begin{esempio}[Operazioni sui caratteri]
\begin{lstlisting}
(char-code #\A)        ; => 65
(code-char 65)         ; => #\A
(char-upcase #\a)      ; => #\A
(char-alphabetic-p #\a) ; => T
(char-digit-p #\5)     ; => 5   (valore numerico, non solo T)
(char= #\a #\a)        ; => T
(char= #\a #\A)        ; => NIL  (case-sensitive)
(char-equal #\a #\A)   ; => T    (case-insensitive)
\end{lstlisting}
\end{esempio}

\argomento{Stringhe}

\begin{concetto}{Il tipo \fnt{string}}
\keyline{\fn{length}}{numero di caratteri.}
\keyline{\fn{char}}{accede all'i-esimo carattere: \fn{(char s i)}.}
\keyline{\fn{string=}, \fn{string<}}{confronto (case-sensitive).}
\keyline{\fn{string-equal}}{confronto case-insensitive.}
\keyline{\fn{string-upcase}/\fn{string-downcase}}{nuova stringa convertita.}
\keyline{\fn{string-trim}}{rimuove caratteri ai margini.}
\keyline{\fn{search}}{posizione di una sottostringa: \fn{(search sub str)}.}

\boxsep
\detail{Le stringhe in Common Lisp sono mutabili (si può usare
\fn{setf} su \fn{char}), ma nella pratica si preferisce
costruire nuove stringhe piuttosto che modificare quelle esistenti.}
\end{concetto}

\begin{esempio}[Operazioni sulle stringhe]
\begin{lstlisting}
(length "ciao")               ; => 4
(char "ciao" 1)               ; => #\i
(string= "abc" "abc")         ; => T
(string= "abc" "ABC")         ; => NIL
(string-equal "abc" "ABC")    ; => T
(string-upcase "ciao mondo")  ; => "CIAO MONDO"
(string-trim " " "  ciao  ")  ; => "ciao"
(search "ao" "ciao mondo")    ; => 2   (indice di inizio)
(search "xyz" "ciao")         ; => NIL (non trovato)

;; Costruire stringhe
(concatenate 'string "foo" "-" "bar")  ; => "foo-bar"
(format nil "~a ha ~a anni" "Mario" 23) ; => "Mario ha 23 anni"
\end{lstlisting}
\end{esempio}

\begin{confronto}
\textbf{Confronto con Java.}
In Java \texttt{String} è immutabile: ogni operazione
crea una nuova stringa. In Common Lisp le stringhe sono
tecnicamente mutabili, ma la convenzione funzionale è
di trattarle come immutabili e costruire nuove stringhe.
\fn{string-upcase} e \fn{string-downcase} restituiscono
una \emph{nuova} stringa senza modificare l'originale,
esattamente come Java.
\end{confronto}

\begin{workbook}
\textbf{Esercizi 137--139}\\[4pt]
L'esercizio~137 (conta vocali) usa \fn{char} e un ciclo.
L'esercizio~138 (palindromo) usa \fn{reverse} su stringhe.
L'esercizio~139 (capitalizza parole) combina
\fn{string-trim}, \fn{search} e \fn{concatenate}.
\end{workbook}

% ================================================================
\section{Stream e file}
% ================================================================

Uno stream è un'astrazione per una sorgente o destinazione
di dati: un file, il terminale, una stringa in memoria.
Le stesse funzioni di lettura e scrittura funzionano
su qualsiasi tipo di stream.

\begin{concetto}{Stream}
\keyline{Input stream}{sorgente: si legge da esso.}
\keyline{Output stream}{destinazione: si scrive su di esso.}
\keyline{\fn{*standard-input*}}{stream standard di input (tastiera/REPL).}
\keyline{\fn{*standard-output*}}{stream standard di output.}
\keyline{\fn{with-open-file}}{apre un file, lega lo stream a una variabile,
chiude automaticamente alla fine.}

\boxsep
\detail{Lo stream si comporta come un cursore: ogni lettura
avanza la posizione. Non si torna indietro senza aprire
il file di nuovo (salvo usare \fn{file-position}).}
\end{concetto}

\begin{esempio}[Leggere un file riga per riga]
\begin{lstlisting}
;; read-line legge una riga come stringa.
;; Il secondo argomento NIL dice "non segnalare errore a fine file".
;; Il terzo argomento NIL e' il valore restituito a fine file.

(defun leggi-righe (percorso)
  (with-open-file (stream percorso :direction :input)
    (loop for riga = (read-line stream nil nil)
          while riga
          collect riga)))

;; Risultato: lista di stringhe, una per riga.
\end{lstlisting}
\end{esempio}

\begin{esempio}[Scrivere su un file]
\begin{lstlisting}
;; :direction :output crea o sovrascrive il file.
;; :if-exists :supersede dice di sovrascrivere se esiste.

(defun scrivi-righe (percorso righe)
  (with-open-file (stream percorso
                          :direction :output
                          :if-exists :supersede
                          :if-does-not-exist :create)
    (dolist (riga righe)
      (write-line riga stream))))
\end{lstlisting}
\end{esempio}

\begin{attenzione}
Attenzione all'encoding. Allegro CL e la maggior parte
delle implementazioni moderne usano UTF-8 per default,
ma il comportamento può dipendere dalla piattaforma.
Per file con caratteri non ASCII, specificare esplicitamente
\fn{:external-format :utf-8} in \fn{with-open-file}.
\end{attenzione}

\begin{workbook}
\textbf{Esercizi 140--141}\\[4pt]
L'esercizio~140 legge un file e conta le righe, le parole
e i caratteri. L'esercizio~141 filtra le righe che
soddisfano un predicato e le scrive su un nuovo file.
\end{workbook}

% ================================================================
\section{\texttt{read} e \texttt{write}: il confine codice/dati}
% ================================================================

Una delle caratteristiche più distinctive di Lisp
è che il suo reader di default capisce la sintassi
del linguaggio stesso.
\fn{read} non parsa stringhe: parsa S-expression.

\begin{concetto}{\fnt{read}, \fnt{write} e il ciclo REPL}
\keyline{\fn{read}}{legge una S-expression da uno stream
(o da \fn{*standard-input*}) e la restituisce come struttura dati Lisp.}
\keyline{\fn{print}, \fn{write}}{serializza una struttura dati
come testo leggibile dal reader.}
\keyline{\fn{read-from-string}}{parsa una stringa come S-expression.}
\keyline{\fn{write-to-string}}{serializza in una stringa.}

\boxsep
\detail{Il ciclo REPL (Read-Eval-Print Loop) è letteralmente:
\fn{(loop (print (eval (read))))}. I tre passi sono funzioni
del linguaggio standard.}
\end{concetto}

\begin{esempio}[Roundtrip struttura dati \(\to\) stringa \(\to\) struttura]
\begin{lstlisting}
;; Serializza una struttura in stringa
(write-to-string '(1 2 (a b) "ciao"))
; => "(1 2 (A B) \"ciao\")"

;; Rilegge la stringa come struttura Lisp
(read-from-string "(1 2 (a b) \"ciao\")")
; => (1 2 (A B) "ciao")

;; Funziona con qualsiasi struttura
(let ((dati (list :nome "Mario" :eta 23 :citta "Roma")))
  (equal dati (read-from-string (write-to-string dati))))
; => T   (roundtrip perfetto)
\end{lstlisting}
\end{esempio}

\begin{attenzione}
\fn{read} esegue i reader macro (\fn{\#'}, \fn{`}, \fn{,}).
Su input non fidato (stringhe da un utente esterno, rete, file
di terze parti) questo è un rischio di sicurezza: un input
malevolo potrebbe far eseguire codice arbitrario al momento
della lettura.

\medskip
Per parsare dati non fidati, usare \fn{read} con
\fn{*read-eval*} impostato a \fn{nil}:
\begin{lstlisting}
(let ((*read-eval* nil))
  (read-from-string stringa-esterna))
\end{lstlisting}
Con \fn{*read-eval*} a \fn{nil}, le reader macro che
eseguono codice (\fn{\#.}) segnalano un errore invece
di eseguirlo.
\end{attenzione}

\begin{workbook}
\textbf{Esercizio 142}\\[4pt]
Implementa un semplice database in memoria che viene
salvato su file con \fn{write} e ricaricato con \fn{read}.
È il caso d'uso pratico più comune di questo meccanismo.
\end{workbook}

% ================================================================
\section{Tokenizzazione}
% ================================================================

La tokenizzazione è il primo passo di qualsiasi parser:
spezzare una stringa in unità significative (token)
prima di interpretarle.

\begin{concetto}{Token e tokenizzazione}
\keyline{Token}{unità minima con significato: numero, simbolo,
operatore, parentesi.}
\keyline{Tokenizzatore}{funzione che prende una stringa e restituisce
la lista dei token.}
\keyline{}{Il tokenizzatore non capisce la struttura: solo le unità.
Il parser costruisce la struttura dai token.}
\end{concetto}

\begin{esempio}[Tokenizzatore su whitespace]
\begin{lstlisting}
;; Spezza una stringa sui caratteri di whitespace.
;; Restituisce una lista di stringhe (i token).

(defun tokenizza (str)
  (let ((token '())
        (corrente '()))
    (map nil
         (lambda (c)
           (if (member c '(#\Space #\Tab #\Newline))
               (when corrente
                 (push (coerce (nreverse corrente) 'string) token)
                 (setf corrente '()))
               (push c corrente)))
         str)
    ;; Flush dell'ultimo token
    (when corrente
      (push (coerce (nreverse corrente) 'string) token))
    (nreverse token)))

(tokenizza "3 4 + 2 *")
; => ("3" "4" "+" "2" "*")
\end{lstlisting}
\end{esempio}

\begin{nota}
\fn{(coerce lista 'string)} converte una lista di caratteri
in stringa. È la funzione opposta di \fn{(coerce stringa 'list)}.
\fn{map nil} è come \fn{mapcar} ma scarta i risultati
e restituisce \fn{nil}: si usa quando l'effetto collaterale
(qui: costruire i token) è il punto.
\end{nota}

\begin{workbook}
\textbf{Esercizio 143}\\[4pt]
Estendi il tokenizzatore per riconoscere non solo stringhe
ma già il tipo di ogni token: \fn{:number} per sequenze
di cifre, \fn{:operator} per \fn{+}, \fn{-}, \fn{*}, \fn{/},
\fn{:symbol} per il resto.
Questo è il passo che rende il parser più semplice.
\end{workbook}

% ================================================================
\section{Parser RPN}
% ================================================================

La notazione polacca inversa (RPN, Reverse Polish Notation)
è il banco di prova ideale per un parser minimale:
la grammatica è semplice, non ci sono parentesi da gestire,
e uno stack è la struttura naturale per l'evaluazione.

\begin{concetto}{RPN e lo stack}
\keyline{}{In RPN gli operatori seguono gli operandi: \texttt{3 4 +} = \texttt{7}.}
\keyline{}{Algoritmo: scorri i token. Se numero: metti sullo stack.
Se operatore: togli due operandi dallo stack, calcola, rimetti il risultato.}
\keyline{}{Alla fine dello stream il risultato è in cima allo stack.}
\end{concetto}

\begin{esempio}[Evaluatore RPN completo]
\begin{lstlisting}
(defun stringa->numero (s)
  "Restituisce il numero se s e' una stringa numerica, NIL altrimenti."
  (let ((*read-eval* nil))
    (ignore-errors
      (let ((v (read-from-string s)))
        (when (numberp v) v)))))

(defun applica-op (op a b)
  (case (intern (string-upcase op))
    (+ (+ a b))
    (- (- a b))
    (* (* a b))
    (/ (/ a b))
    (t (error "Operatore sconosciuto: ~a" op))))

(defun valuta-rpn (espressione)
  "Valuta una stringa in notazione RPN.
   Es: (valuta-rpn \"3 4 + 2 *\") => 14"
  (let ((stack '()))
    (dolist (token (tokenizza espressione))
      (let ((n (stringa->numero token)))
        (if n
            (push n stack)
            (let ((b (pop stack))
                  (a (pop stack)))
              (push (applica-op token a b) stack)))))
    (first stack)))

(valuta-rpn "3 4 +")        ; => 7
(valuta-rpn "3 4 + 2 *")    ; => 14
(valuta-rpn "5 1 2 + 4 * + 3 -") ; => 14
\end{lstlisting}
\end{esempio}

\begin{nota}
\fn{ignore-errors} esegue il corpo e, se viene segnalata
una condizione, restituisce \fn{nil} invece di propagare
l'errore. Utile per conversioni tentative (``prova a leggere
un numero, se non riesci restituisci nil'').
Per una gestione più robusta degli errori vedi il Capitolo~9.
\end{nota}

\begin{workbook}
\textbf{Esercizio 144}\\[4pt]
Estendi \fn{valuta-rpn} con le operazioni \fn{mod}, \fn{expt}
e \fn{sqrt}. Poi aggiungi una gestione degli errori che
segnali un messaggio leggibile quando lo stack è vuoto
o quando si incontra un token non riconosciuto.
\end{workbook}

% ================================================================
\section{Mini-evaluator per S-expression}
% ================================================================

Il passo finale: valutare un sottoinsieme di S-expression Lisp.
Non è un esercizio accademico: è esattamente come funziona
il cuore di un interprete Lisp.

\begin{intuizione}
Una S-expression è già una struttura dati Lisp.
\fn{read-from-string "(+ 1 (* 2 3))"} restituisce
\fn{(+ 1 (* 2 3))}: una lista con testa \fn{+} e argomenti
\fn{1} e \fn{(* 2 3)}.
Valutarla significa: riconosci il tipo di espressione
(numero, simbolo, lista), poi applica la regola giusta.
\end{intuizione}

\begin{concetto}{La struttura di un evaluator}
\keyline{Numero}{si valuta a se stesso.}
\keyline{Simbolo}{cerca nell'ambiente.}
\keyline{Lista \fn{(f a1 a2...)}}{valuta \fn{f} per ottenere la funzione,
valuta ricorsivamente ogni argomento, applica.}
\keyline{Ambiente}{tabella che associa nomi a valori (hash table o alist).}
\end{concetto}

\begin{esempio}[Mini-evaluator su un sottoinsieme di Lisp]
\begin{lstlisting}
;; Ambiente di default: le operazioni aritmetiche.
(defparameter *env-default*
  (list (cons '+ #'+)
        (cons '- #'-)
        (cons '* #'*)
        (cons '/ #'/)
        (cons 'abs #'abs)
        (cons 'max #'max)
        (cons 'min #'min)))

(defun cerca-env (simbolo env)
  (let ((coppia (assoc simbolo env)))
    (if coppia
        (cdr coppia)
        (error "Simbolo non definito: ~a" simbolo))))

(defun mini-eval (expr &optional (env *env-default*))
  (cond
    ;; Caso base 1: numero (o stringa) -> se stesso
    ((numberp expr) expr)
    ((stringp expr) expr)
    ;; Caso base 2: simbolo -> cerca nell'ambiente
    ((symbolp expr) (cerca-env expr env))
    ;; Caso ricorsivo: lista (f a1 a2 ...)
    ((listp expr)
     (let ((fn   (mini-eval (first expr) env))
           (args (mapcar (lambda (a) (mini-eval a env))
                         (rest expr))))
       (apply fn args)))
    (t (error "Espressione non riconosciuta: ~a" expr))))

;; Uso
(mini-eval '(+ 1 2))             ; => 3
(mini-eval '(* 2 (+ 3 4)))       ; => 14
(mini-eval '(max 1 5 3 2))       ; => 5
(mini-eval '(abs (- 3 7)))       ; => 4

;; Con read-from-string: valuta stringhe direttamente
(let ((*read-eval* nil))
  (mini-eval (read-from-string "(+ 1 (* 2 3))")))
; => 7
\end{lstlisting}
\end{esempio}

\begin{intuizione}
Questo è il cuore di Lisp.
Il REPL di qualsiasi implementazione Common Lisp
fa esattamente questo: \fn{read} parsa, \fn{eval} valuta
ricorsivamente seguendo le stesse regole.
La differenza tra \fn{mini-eval} e il vero \fn{eval} di Lisp
è nei casi gestiti: il vero \fn{eval} gestisce anche
\fn{if}, \fn{let}, \fn{lambda}, \fn{quote}, le macro, ecc.
Ma la struttura è identica.
\end{intuizione}

\begin{esempio}[Estendere l'ambiente con variabili]
\begin{lstlisting}
;; Aggiungere variabili all'ambiente: basta fare acons.
(let ((env (acons 'x 10 (acons 'y 3 *env-default*))))
  (mini-eval '(+ x (* y 2)) env))
; => 16

;; Notare: l'ambiente e' una alist.
;; acons aggiunge in testa senza modificare *env-default*.
;; La variabile in testa fa "shadow" su eventuali omonimi.
\end{lstlisting}
\end{esempio}

\begin{workbook}
\textbf{Esercizi 145--148}\\[4pt]
L'esercizio~145 aggiunge \fn{if} al mini-evaluator (con
valutazione condizionale degli argomenti).
L'esercizio~146 aggiunge \fn{let} (estende l'ambiente localmente).
L'esercizio~147 aggiunge \fn{lambda} (chiusura sull'ambiente corrente).
L'esercizio~148 è il progetto finale: un mini-REPL
che legge espressioni da terminale, le valuta con
\fn{mini-eval} e stampa il risultato.\\[4pt]
L'esercizio~148 è la sintesi di tutti i capitoli precedenti:
strutture dati (alist per l'ambiente), higher-order
(funzioni come valori nell'ambiente), ricorsione
(la valutazione è ricorsiva), I/O (legge da terminale).
\end{workbook}

% ================================================================
\section*{Riepilogo del capitolo}
% ================================================================

\begin{tcolorbox}[
  enhanced,
  colback=csBlue!5, colframe=csBlue!40,
  arc=4pt, boxrule=0.6pt,
  left=14pt, right=14pt, top=12pt, bottom=12pt,
  before skip=16pt, after skip=16pt
]
\textbf{Quello che hai imparato:}

\medskip
\begin{itemize}
  \item \textbf{Sequence}: tipo astratto comune a liste, vettori
        e stringhe. \fn{map}, \fn{subseq}, \fn{concatenate}
        richiedono il tipo del risultato come primo argomento.
  \item \textbf{Caratteri e stringhe}: caratteri con \fn{\#\textbackslash{}},
        stringhe come vettori di caratteri. \fn{char-code}/\fn{code-char},
        \fn{string=}, \fn{string-trim}, \fn{search}.
  \item \textbf{Stream e file}: \fn{with-open-file} per la lettura
        e scrittura. Lo stream è un cursore: avanza a ogni lettura.
  \item \textbf{\fn{read} e \fn{write}}: confine tra codice e dati.
        \fn{read-from-string} e \fn{write-to-string} per il roundtrip.
        Mai usare \fn{read} su input non fidato senza
        \fn{*read-eval* nil}.
  \item \textbf{Tokenizzazione}: prima fase di ogni parser.
        Spezza la stringa in unità minimali prima di interpretarla.
  \item \textbf{RPN}: stack + un passaggio sui token.
        Pattern trasferibile a qualsiasi grammatica context-free semplice.
  \item \textbf{Mini-evaluator}: numero, simbolo, lista.
        La ricorsione sulla struttura della S-expression
        è il cuore di ogni interprete Lisp.
  \item \textbf{Ambiente come alist}: \fn{acons} aggiunge senza mutare,
        \fn{assoc} cerca. Lo stesso meccanismo del Capitolo~5,
        ora usato come ambiente di valutazione.
\end{itemize}

\medskip
\textbf{Prossimo:} Capitolo~9, \emph{Stile e professione}.
Come organizzare il codice, documentarlo, testarlo e portarlo
in produzione.
\end{tcolorbox}

\begin{workbook}
\textbf{Prima di andare al Capitolo~9:}\\[4pt]
Completa gli esercizi \textbf{136--147} (core del capitolo).
L'esercizio~148 (mini-REPL) è il progetto conclusivo:
affrontalo dopo aver fatto gli altri.\\[4pt]
Priorità: \textbf{137, 140, 142, 144, 145, 148}.
Questi sei coprono stringhe, file, \fn{read}/\fn{write},
RPN, evaluator con \fn{if} e il progetto finale.
\end{workbook}
